\clearpage\section{Overall architecture}\label{sec:architecture}
\subsection{Container and deployment view}
Figure~\ref{fig:containers} shows the proposed GENAILMS services and data stores. The browser runs React and TypeScript. Caddy routes application requests. FastAPI serves the API and starts a background worker thread in the backend process. The queue resides in PostgreSQL. Moodle and its bridge plugin form the external delivery boundary. Moodle is reached only through its Web Services and the \code{local_lms_bridge} plugin; the application PostgreSQL store belongs to the authoring system.
\diagram{containers}{Container/deployment architecture and service calls.}{containers}
\subsection{Service responsibilities}
\par
Table~\ref{tab:04-architecture-1} presents service boundaries and persistent state.
\begin{longtable}{|p{2.7cm}|p{6.0cm}|p{5.2cm}|}\caption{Service boundaries and persistent state.}\label{tab:04-architecture-1}\\\hline
Service & Responsibility & Data / boundary\\\hline\endfirsthead
\hline Service & Responsibility & Data / boundary\\\hline\endhead
Frontend & Routes, authoring forms, polling, preview, editing, blind-rating form & Browser session; API-derived content\\\hline
Backend API & Authentication, validation, ownership, generation requests, editing and publication & Runtime DB role; request-scoped user context\\\hline
Backend worker & Claim jobs, execute ingestion/generation, classify errors and retry & Worker DB role; cross-owner processing with restricted grants\\\hline
PostgreSQL & Courses, documents, chunks, modules, jobs, sessions, studies and ratings & Transactions, foreign keys, RLS, pgvector\\\hline
Upload storage & Keep files until extraction/embedding completes & Generated filenames; staged files deleted after processing\\\hline
LLM gateway & Provider keys, alias mapping, structured calls, retries and fallback & Configuration and provider boundary\\\hline
Embedder & Local Qwen3 document/query vectors & Model cache and pinned Open Neural Network Exchange (ONNX) revision\\\hline
Reranker & Cross-encoder ranking of candidate passages & Local Jina multilingual INT8 model cache\\\hline
Moodle bridge & Identity confirmation, actor permissions and Book/Quiz writes & Moodle-owned identity, context and activities\\\hline\end{longtable}
\subsection{Why this architecture}
A shared PostgreSQL store with pgvector\citep{pgvector} avoids introducing another stateful service for a course-sized workload. Queue claiming uses a row lock with \code{SKIP LOCKED}\citep{pgselect}, allowing a claimant to skip work already locked elsewhere. Vector search uses exact cosine ranking within a filtered course/document subset. Course/document filtering precedes vector ranking. This is a scale-specific choice: an increasing number of chunks or concurrent jobs would justify new measurement and possibly a different indexing or worker design.

Extraction, embedding, reranking and external generation have different runtime profiles. Local models run in dedicated containers, keeping model-serving dependencies out of the backend image. Generation calls pass through a single provider gateway, so provider credentials and routing changes do not spread through every pipeline function. Embeddings traverse the gateway to the local embedder; reranking is a direct backend-to-reranker request.
\subsection{Transaction and failure boundaries}
An API transaction records the work request; the asynchronous worker claims and saves a job in a separate transaction. Ingestion records step transitions so a failed embed step need not re-extract the file. Before slow embedding, the worker ends its read transaction. Retrieval for generation uses a separate reading session before reranking to avoid holding that transaction during model inference.

Generation has a coarser boundary. Its saved job step invokes the module-generation chain; failure can rerun the chain and incur additional model calls. A cached planner reduces repeated planning but is not a checkpoint for each generated part. Publication is another boundary: local database rollback cannot undo already created Moodle activities. Phase 2 will evaluate retries, resource use and publication recovery.
